\chapter{有限诊断中的思维过程}
\label{ch:model}
\chapterguide{先固定系统能够观察什么，再比较它怎样组织观察。未知情境的真值不进入控制者输入。}

\section{一个持续的共同错误}
设目标为事实$Y\in\{0,1\}$，某条推理规则受到隐藏偏差$B\in\{0,1\}$影响。重复使用它可能稳定得到$Y\oplus B$。多数表决可以压低独立噪声，却无法决定稳定答案来自真实事实还是共同偏差。

每次调用的扰动会重新发生；偏差$B$则贯穿整段运行。用固定但未知的情境$\theta$同时记录真实目标与持续因素。不同成员即使使用独立随机种子，也可能共同受到$B$影响。

\section{片段、操作与推理链模式}
片段$v$可以是一项观察、一个候选解释或已经组合好的记录。一次推理写作
\[
 \II_e\xrightarrow{F_a}v,\qquad \II_e=[v_{i_1},\ldots,v_{i_r}].
\]
有序输入有向超边代表一次操作。两次独立调用分别产生两个节点，同一节点可以被多个操作读取。

推理链$\LL_t=[e_1,\ldots,e_t]$包含操作及其输入输出。推理链模式$\Pi$依据这条记录和外部到达的可用材料，选择下一次绑定了输入的操作、等待或停止。执行模式并持续保有自身历史与可见材料的运行实例称为Agent。模式可以由另一Agent提供；生成它所用的操作也进入总成本。

片段层自回归模式规定下一输入追加刚得到的片段；其他模式可以重新选择和排列片段。单次$F_a$内部是否逐token生成，属于另一层描述。

多Agent组织的成员根据各自局部历史行动。证明下界时，可以让更强的比较者看到全部记录；构造分布式算法时，必须说明记录通过什么消息到达所需位置。

\section{受控的有限实验族}
\begin{definition}[诊断实例]
实例由有限集合$\Theta,\calA,\VV$、目标函数$g:\Theta\to\mathcal Y$和概率表$K_a(v\mid\theta)$组成。第$t$次调用满足
\begin{equation}
 \PP_\theta(V_t=v\mid T_{t-1},A_t=a)=K_a(v\mid\theta),
 \quad T_t=[(A_1,V_1),\ldots,(A_t,V_t)].
 \label{eq:controlled}
\end{equation}
控制者知道概率表与目标函数，未知真实$\theta$，只能通过调用取得有关$\theta$的新信息。
\end{definition}

给定$\theta$与所选实验，下一观察不再依赖更早记录；选择哪个实验仍可依赖全部历史。持续错误通过$\theta$进入各次调用。对未知情境混合后，观察一般相关。

实验$a$是一种允许的查询或处理方式，有限输出$v$就是形成的片段。这一专门化模型研究证据取得与组织。历史会改变基本操作分布的系统，需要增添显式状态；第\ref{ch:frontiers}章讨论新的困难。

\section{完整组织类别与成本}
允许所有依据历史选择实验的随机策略，其内部随机性与真实情境独立。结束时输出$\widehat Y$。定义
\[
 R_n^*=\inf_{\Pi:\,N\le n\ \mathrm{a.s.}}
 \max_\theta\PP_\theta^\Pi(\widehat Y\ne g(\theta)),
 \qquad n^*(\delta)=\min\{n:R_n^*\le\delta\}.
\]
集合为空时$n^*(\delta)=\infty$。停止后用不产生信息的空操作补齐记录，便能同时分析提前停止策略。

主定理计数等成本基本调用。控制者的查表、计数、似然计算和通信另行报告。下界允许控制者使用任意计算，因此也约束受限控制者；上界给出实现所需运算。若调用成本为$c_a$，应增加$\sum_ac_aN_a\le B$约束，并用同样权重改写覆盖规划。

统计定理使用已知概率核。实现部分以显式有理概率表为输入，对平方根、对数和分配作有误差控制的近似，再核对实际覆盖下界。精度成本另计，不把任意实数视为免费可读的无限信息。第\ref{ch:representation}章分析这类数值误差如何进入风险。

\section{一个可以核对的运行}
取四个情境$(Y,B)\in\{0,1\}^2$，两个实验分别观察$Y\oplus B$与$B$，各以概率$\eta<1/2$翻转。策略先重复第一种实验，再决定是否校准第二种。最终输出从观察中得到，控制者不能直接读取$Y$或$B$。

这里没有预设永远可靠的检查器。操作能否区分需要不同答案的情境，将从概率表推出。
\begin{exercise}把全部随机性写入一个预先采样的变量，为什么不意味着控制者能够读取这个变量的真实值？\end{exercise}
\begin{exercise}构造两个拥有相同最后片段、但适合选择不同下一操作的历史，解释模式为何需要使用运行记录。\end{exercise}
